\section{Conclusion}\label{sec:conclusion}

\subsection{Summary}

This paper presented the Seka template, a comprehensive \LaTeX\ framework for academic writing. The template provides:

\begin{itemize}
    \item A professional two-column layout conforming to journal standards
    \item Modular section organization for complex documents
    \item Custom environments for notes, important information, and code
    \item Robust cross-referencing and citation management
    \item No external dependencies or special font requirements
\end{itemize}

\subsection{Key Contributions}

The primary contributions of this work are:

\begin{enumerate}
    \item A ready-to-use academic template that compiles without errors
    \item Demonstration of best practices for \LaTeX\ document organization
    \item A customizable framework adaptable to various publication venues
\end{enumerate}

\subsection{Recommendations}

Based on our experience with this template, we recommend that users:

\begin{itemize}
    \item Maintain the modular structure for all substantial documents
    \item Use the provided custom environments for consistent formatting
    \item Extend the class file rather than modifying individual documents
\end{itemize}

\subsection{Future Directions}

Future work will focus on expanding the template's capabilities while maintaining its simplicity. Specific directions include:

\begin{itemize}
    \item Development of a one-column variant
    \item Integration with reference managers
    \item Enhanced support for supplementary materials
\end{itemize}

\subsection{Final Remarks}

The Seka template is freely available for academic use. Researchers are encouraged to adapt it to their specific needs and contribute improvements back to the community.